% Transcription — fonds Alexandre Grothendieck, shelfmark 64, batch 4 (pages 61-80).
% Established from the facsimile archives/batches/64/batch-04.pdf (a copy of 64.pdf
% fetched through the project relay, no cover sheet: PDF page k = archivists' page 20(N-1)+k),
% per the skill transcribe-grothendieck. The folder is 156 pages in eight batches.
%
% Pass: Opus 5.5 (claude-opus-5-5), 2026-09-27 — first pass, unchecked against the pages by a human.
%
% Pages skipped: 69 (a coloured sheet bearing only his title « III. Épinglages de Legendre
%   des courbes elliptiques », used as the section heading, with a note).
% His pagination: first run continues — p. 62 = his 10, 64 = 11, 66 = (12), 68 = 13
%   (even archive pages only). A new run starts at p. 74 = his 1; 75 = 2, 77 = 3, 79 = 4
%   (his numbers on 76, 78, 80 not visible); its formula numbering restarts at (1).
% What the batch is about: end of the level-2 analysis (the map f_2^t on Sigma_J induced by
%   2 id_{E_0}, computed in Legendre coordinates: f_2^t(z) = t((z^2-2z+t)/(z^2-2tz+t))^2);
%   section 4 « Détermination de M_11(4) » (level-4 structure M = 4E as an extension of
%   M_0 = 2E, V_J(A) = Ker(A^J -> A), a torsor under End(M_0)); then part III, Legendre
%   pinnings: a digression on tamely ramified Kummer covers pinned by tangent spaces
%   (T_{X'}^{⊗n} ≅ T_X), then E described by (J, t, T, T^{⊗2} ≅ T_{Σ,t}), the twist of O_t(1)
%   by a mu_4-torsor Q_0 (a combinatorial square), Legendre and Legendre–Jacobi pinnings,
%   equivalence with pairs (t, i), t in U_{0,3}, i^2 = -1, and the relative version over
%   Spec Z[1/2] with the twist M^natural = M ∧_{mu_4} Q_0.
% Notation fixed once for the batch: \underline{J} (étale cubic cover), \Sigma_J, D_J,
%   \Sigma_J^*, \underline{\omega} (orientations), \mu_4^* = {i,-i}, Q_0 (square), D_{Q_0}.
%   His zigzag twist superscript (pp. 79-80) is rendered $\natural$; his b-like
%   superscript in (23) as $\flat$.
% \ill{}: 67. \uncertain{}: 42.

\documentclass[11pt,a4paper]{article}
\input{../preamble/grothendieck}

\folder{64}
\batch{4}
\pages{61}{80}
\dating{[à partir de 1980- 1982]}
\watermark{Édition de démonstration}
\foldertitle{Modules courbes elliptiques en caractéristiques ≠2 : notes manuscrites (s.d.)}

\begin{document}

\page{61}
(au dessus duquel $X^*$ est étale \struck{galoi} principal de
groupe $\underline{V}\}$, tel comme $E$ sur $E_0$) mais non
\ill{} ramifiés aux pts de $D_J \subset \Sigma_J$
(vu que $E_0/\Sigma_J$ y est ramifié, mais non $E/X$ au-dessus des points
de $\Delta = X|D_J$, or si c'était contraire il en résulterait que
$E_0/\Sigma_J$ est étale \struck{sur} sur les pts de $D_J$ aussi, par
descente plate\ldots). (Quand \ill{} $=$ il doit être facile de vérifier
que le diagramme total se \uncertain{reconstitue} : \uncertain{partir} de
$E_0$ ($\longmapsto \Sigma = \Sigma_{\underline{J}}$, \add{et $\underline{J}$,
et partant de $\Sigma^*$}), et $E_0 \to \Sigma$,) et des revêtements
\struck{principal} $X$ de $\Sigma$ \add{\uncertain{étales}} avec groupe
$\underline{V}$ : la donnée de $(E_0, I)$ \struck{doit être} \add{est}
équivalente \ill{} $(E, \sigma)$ i.e. de \uncertain{couples} : $E_0$ et des
revêtements $X$ (d'où un \uncertain{faisceau} $\underline{K}$ et
$\underline{I} \simeq \underline{K}_1(X|t)$\ldots)

Cela montre que \uncertain{localement} $(E, \sigma)$ i.e. le $I$ est un
\ill{} de la \ill{} \ill{}, on se ramène par le cas $E = E_0$,
$\sigma = \sigma_0$, donc $I = \underline{V} = {}_2E_0 = M_0$, torseur
trivial. On a donc \struck{\ill{}} $X = \Sigma_J$ et on trouve donc

\page{62}\note{p.~10 de l'auteur, chiffrée en tête de page.}
une application canonique
\[
  (30) \qquad \Sigma_{\underline{J}} \xrightarrow{\;f_2^t\;} \Sigma_{\underline{J}}
\]
\marginal{c'est l'application déduite par passage au quotient de
$2\,\mathrm{id}_{E_0} : E_0 \to E_0$}
qui doit pouvoir se définir en termes de $\underline{J}$, et de la seule
section $t$ de $\Sigma_{\underline{J}}$ (ce qui me semble déterminer $E_0$,
à isom\supplied{orphismes} (non uniques) près loc.\ pour la top.\ étale\ldots)
On a ici
\struck{$\underline{I} = \underline{J} \amalg$}
\[
  (31) \qquad \underline{I}_t = \underline{J} \amalg t(S)
\]
\marginal{\emph{NB} comme $S$-schéma, $\underline{I}_t$ ne dépend pas de
$t$, mod\supplied{ulo} iso.\ can.}
\marginal{il faudrait revenir aux \ill{} descriptions intrinsèques \ill{}
des \ill{} \uncertain{intervenant} \ill{} courbes elliptiques}
\note{la seconde note marginale est écrite en diagonale dans la marge gauche,
à hauteur de (30), et reliée par un trait double au « qui doit pouvoir se
définir » ; elle n'est lue qu'en partie.}
comme $V_{\underline{I}} \simeq V_{\underline{J}}$, et \struck{$E_0$}
$\underline{V}$ peut se décrire comme opérant sur $\underline{I}$ de la
façon canonique, \struck{cette} opération se prolongeant en une opération
sur $\Sigma_{\underline{J}}$ (dépendant de la section $t$ de
$\Sigma_{\underline{J}}^*$). L'\struck{application} morphisme $f_t$ est
défini sans ambiguïté par les conditions suivantes (où on a
\uncertain{supposé} pour simplifier que $\underline{J} = J_S$, $J$ un
\ill{} : trois éléments) :
\[
  (32) \quad
  \begin{cases}
    \forall\, i \in J, \text{ si } \sigma_i^t \in \operatorname{Aut}_S(\Sigma_J)
      \text{ est l'automorphisme qui échange } t \text{ et } i, \\
    \qquad \text{et } j \text{ et } k \ (\text{où } \{i,j,k\} = J),
      \text{ et } \Delta_i^t = (\Sigma_J)^{\sigma_i^t}, \\
    f(\Delta_i^t) = \text{section } \xi_i \text{ de } D_{\underline{J}} \simeq J_S
  \end{cases}
\]
\note{avant la seconde ligne de (32), un « alors $f(\Sigma_J^{\sigma_i})$ »
biffé.}
et cela doit impliquer automatiquement
\[
  (33) \qquad f_t\bigl(I_S = J_S \cup t(S)\bigr) \subset t(S)
\]
\note{devant $t(S)$, le mot « section » est biffé.}

\page{63}
Faisons les calculs pour $J = \{0, 1, \infty\} \subset \mathbb{P}^1_S$,
on trouve facilement :
\[
  (34) \quad
  \begin{cases}
    \sigma_0^t(z) = \dfrac{z - t}{z - 1} & (\text{éch.\ } 0 \text{ et } t,\ 1 \text{ et } \infty) \\[1ex]
    \sigma_1^t(z) = t/z & (\text{éch.\ } 1 \text{ et } t,\ 0 \text{ et } \infty) \\[1ex]
    \sigma_\infty^t(z) = t\,\dfrac{z - 1}{z - t} & (\text{éch.\ } \infty \text{ et } t,\ 0 \text{ et } 1)
  \end{cases}
  \qquad \text{NB } \sigma_\infty^t = \sigma_0^t \sigma_1^t
\]
On trouve donc
\[
  (35) \quad
  \begin{cases}
    \Delta_0^t : \text{solutions de } z^2 - 2z + t = 0, & \Delta_0^t = \{1 \pm \sqrt{1-t}\} \\
    \Delta_1^t : \text{\phantom{solutions de }} z^2 - t = 0 & \Delta_1^t = \{\pm \sqrt{t}\} \\
    \Delta_\infty^t : \text{\phantom{solutions de }} z^2 - 2tz + t = 0 & \Delta_\infty^t = \{t \pm \sqrt{t(1-t)}\}
  \end{cases}
\]
\note{les deuxième et troisième lignes répètent « solutions de » par un trait.}
et on doit avoir
\[
  (36) \qquad f_2^t : \quad
  \begin{aligned}
    \Delta_0^t &\longrightarrow 0 \\
    \Delta_1^t &\longrightarrow 1 \\
    \Delta_\infty^t &\longrightarrow \infty
  \end{aligned}
  \qquad \underline{\text{avec multiplicité } 2}
\]
ce qui nous donne (puisque les zéros et pôles de $f_2^t$ de degré quatre
sont connus)
\[
  f_2^t(z) = c_t \left( \frac{z^2 - 2z + t}{z^2 - 2tz + t} \right)^2
  \qquad c_t \text{ section de } \underline{\mathcal{O}}_S^*,
\]
dont la valeur s'obtient en écrivant que $\Delta_1^t \to 1$. Mais si $\zeta$
est une solution de $\Delta_1^t$ i.e.\ $\zeta^2 = t$, on a
\[
  \begin{aligned}
    \zeta^2 - 2\zeta + t &= 2(t - \zeta) \\
    \zeta^2 - 2t\zeta + t &= 2t(1 - \zeta)
  \end{aligned}
\]
\[
  f_2^t(\zeta) = c_t \left( \frac{2(t - \zeta)}{2t(1 - \zeta)} \right)^2
  = c_t/t^2 \left( \frac{t - \zeta}{1 - \zeta} \right)^2
\]
\note{entre $c_t/t^2$ et le dernier facteur, une fraction biffée,
\struck{$\frac{\zeta^2 - 2t\zeta + t^2}{\zeta^2 - 2\zeta + 1}$}, remplacée
par $\left(\frac{t-\zeta}{1-\zeta}\right)^2$.}
Or \struck{$\frac{t^2 + t - 2t\zeta}{1 + t - 2\zeta}$}
\[
  \frac{t - \zeta}{1 - \zeta}
  = \frac{(t - \zeta)(1 + \zeta)}{(1 - \zeta)(1 + \zeta)}
  = \frac{-\zeta^2 + (t-1)\zeta + t}{1 - \zeta^2}
  = \frac{(t-1)\zeta}{1 - t} = -\zeta
\]
donc
\[
  f_2^t(\zeta) = \frac{c_t}{t^2}\, \zeta^2 = \frac{c_t\, t}{t^2} = \frac{c_t}{t}
  \qquad \text{d'où} \qquad c_t = t
\]
\note{les deux calculs biffés sont barrés de traits obliques ; la
première ligne de l'identité en $\zeta$ est surchargée (un « $t$ » corrigé
devant $-2\zeta$).}

\page{64}\note{p.~11 de l'auteur, chiffrée en tête de page.}
et par suite
\[
  (37) \qquad f_2^t(z) = t \left( \frac{z^2 - 2z + t}{z^2 - 2tz + t} \right)^2
\]
Il reste : vérifier que les valeurs de cette fonction pour
$z = 0, 1, \infty$ et $t$, sont égales à $t$, i.e.\ que celle de
$\dfrac{z^2 - 2z + t}{z^2 - 2tz + t}$ est égale à $\pm 1$. Pour
$0, \infty, 1$ on trouve $1$, pour $z = t$ on trouve
$\dfrac{t^2 - t}{t - t^2} = -1$, OK.

\emph{NB} En prenant la $n$-ième itérée de $f_2^t$, on trouve le
morphisme $f_{2^n}^t$, qui correspond par passage au quotient à
$2^n \mathrm{id}_{E_0}$. L'image inverse de la section $t$ par
$f_{2^n}^t = (f_2^t)^{\circ n}$ est l'image par
$p : E_0 \to \Sigma_J \simeq E_0/\{1, \sigma_{E_0}\}$ du \struck{l'ensemble}
\add{sous-schéma en groupes} ${}_{2^n}E_0$ des points annulés par $2^n$.

\section{4. Détermination de $M_{11}(4)$.}

\page{65}
Soit $E$ une courbe elliptique sur $S$ (où car.\ résiduelles $\neq 2$),
considérons
\[
  (38) \qquad M = {}_4E
\]
schéma en groupes sur $S$ loc.\ isom.\ à $(\mathbb{Z}/4\mathbb{Z})^2$, i.e.\
\struck{syst.\ local} faisceau abélien \struck{pour} sur $S_{\text{ét}}$ qui
est un $(\mathbb{Z}/4\mathbb{Z})_S$-module loc.\ libre de rg $2$.

\struck{(39)} On sait qu'on a


\begin{tikzcd}
  0 \arrow[r] & M_0 \arrow[r] & M \arrow[r, "\mathrm{can}"] & M_0 \arrow[r] & 0 \\
  & M \arrow[u] \arrow[ur, "2\,\mathrm{id}_M"'] & & &
\end{tikzcd}

(39), où
\[
  M_0 = M \otimes_{\mathbb{Z}/4\mathbb{Z}} \mathbb{F}_2 \simeq {}_2E .
\]
\note{le numéro (39) est écrit en marge gauche de la suite exacte, un premier
« (39) » biffé au-dessus ; la flèche verticale $M \to M_0$ est tracée vers le
haut, sans étiquette lisible.}

La donnée d'un $M$ équivaut à celle d'un $M_0$ (i.e.\ d'un $\underline{J}$
\struck{rev.} rev.\ étale de degré $3$ sur $S$ \struck{\ill{}}
\[
  (40) \qquad \underline{J} \simeq M_0^*, \quad
  M_0 = V_{\underline{J}}(\mathbb{F}_2) \simeq
  \mathbb{F}_2^{\underline{J}} / \underset{\text{diag}}{\mathbb{F}_2}\ ),
\]
\emph{et} d'un torseur \add{$P$} sous
\[
  (41) \qquad \text{« } V_{\underline{J}}(M_0) \text{ »} \simeq
  \underline{\mathrm{End}}_{\mathbb{F}_2, S}(M_0) \simeq
  M_0 \oplus \mathbb{F}_2^{\underline{\omega}}
\]
\marginal{$\underline{\omega}$ est le rev.\ étale de degré $2$ « des
orientations » associé au rev.\ $\underline{J}$ de degré $3$ de $S$}
Un épinglage de $M = {}_4E$ (i.e.\ \struck{d'un} la donnée d'un iso
\[
  (42) \qquad (\mathbb{Z}/4\mathbb{Z})^2_S \simeq M \ (= {}_4E) \ )
\]
équivaut à : celle d'une \add{trivialisation de $\underline{J}$, plus d'une}
section (trivialis.\ \ill{}
\note{la note marginale est écrite de biais dans la marge gauche, à hauteur
de (41) ; « des orientations » y est lu d'après le contexte. La dernière
ligne se poursuit page~66.}

\page{66}\note{p.~12 de l'auteur, chiffrée « (12) » en tête de page.}
\struck{plus celle d'une} \add{du torseur}, de ce torseur,
\struck{qui} (peut aussi s'interpréter ainsi : la structure d'extension (39)
\struck{comme le torseur des sections} définit la structure d'extension
\[
  (43) \qquad 0 \to
  \overset{\underline{\mathrm{End}}(M_0)}{\overset{\wr}{V_{\underline{J}}(M_0)}}
  \longrightarrow V_{\underline{J}}(M) \longrightarrow
  \overset{\underline{\mathrm{End}}(M_0)}{\overset{\wr}{V_{\underline{J}}(M_0)}}
  \longrightarrow 1
\]
où pour $A$ schéma \add{(ou faisceau)} en groupes commut.\ $A$ sur $S$, on
pose
\[
  (44) \qquad V_{\underline{J}}(A) \simeq
  \operatorname{Ker}\bigl( A^{\underline{J}} \xrightarrow{\ \text{somme}\ } A \bigr)
\]
(faisceau loc.\ isom.\ à $A^2$), et $V_{\underline{J}}$ est un foncteur exact
en $A$, commutant \ill{} bien sûr aux chgt.\ de base. Pour $A$ annulé par
$2$, on a un isom
\[
  (45) \qquad V_{\underline{J}}(A) \simeq \operatorname{Hom}(M_0, A)
  \qquad \text{si } 2\,\mathrm{id}_A = 0
\]
dû au fait que l'on a
\[
  (46) \qquad M_0 \simeq \mathbb{F}_2^{\underline{J}} / \underset{\text{diag}}{\mathbb{F}_2}
\]
\note{dans (43), l'auteur écrit « $\underline{\mathrm{End}}(M_0)$ » au-dessus
des deux termes extrêmes, reliés à eux par un signe $\wr$ (isomorphisme
vertical).}

Donc à une section \struck{de} canonique partout non nulle de
$V_{\underline{J}}(M_0)$, correspondant à l'application d'inclusion
$\underline{J} \simeq M_0^* \hookrightarrow M_0$, ou encore : la section
évidente $\mathrm{id}_{M_0}$ de $\underline{\mathrm{End}}(M_0)$, dont l'image
inverse dans (43) est le torseur cherché.

\page{67}
\struck{Un épinglage de $M$ s'identifie donc à un}

On peut considérer, indépendamment de la question d'épingler $J$,
\struck{celle} \add{l'épinglage} \struck{\ill{}} « relatif »
\struck{consistant en sections} \add{consistant} en \struck{un} la donnée
d'une section du torseur $P_M$ \add{sur} sous
$\underline{\mathrm{End}}(M_0)$, ce qui équivaut aussi à la donnée d'un
morphisme
\[
  (47) \qquad \underline{J} \longrightarrow M \,\}
\]
\note{devant $M$, un symbole noirci (sans doute $V_{\underline{J}}($) ; l'accolade
fermante est de l'auteur.}
qui relève le morphisme d'inclusion
\[
  \underline{J} \hookrightarrow M_0
\]
et qui soit une section de $V_{\underline{J}}(M)$ i.e.\ qui soit à
\uncertain{trace} nulle (\ill{} : $S$). Ceci en fait automatiquement un
morphisme
\[
  \underline{J} \longrightarrow M \setminus M_0
\]
donc en composant avec \struck{l'inclusion}
\[
  (\underline{J} \dashrightarrow)\; M \hookrightarrow E \longrightarrow
  \Sigma_{\underline{J}} \simeq E/(1, \sigma_E)
\]
on trouve un morphisme
\[
  (48) \qquad \underline{J} \longrightarrow \Delta_t
\]
et la condition de relèvement pour (47) signifie aussi que le composé de (48)

\page{68}\note{p.~13 de l'auteur, chiffrée en tête de page.}
avec le morphisme canonique (étale de degré $2$)
\[
  f_t : \Delta_t \longrightarrow \underline{J}
\]
soit l'identité --- i.e.\ (48) définit une \emph{face} du fibré en
octaèdres combinatoires $\Delta_t$ sur $S$. Du coup l'octaèdre est
$\underline{J}$-épinglé, en d'autres termes, il est can.\ isomorphe au
\struck{\ill{}} octaèdre standard défini par $\underline{J}$, \add{et une
section de $\mu_4^*$,} en prenant les revêtements biquadratiques
$\Sigma_{\underline{J},2}$ de $\Sigma_{\underline{J}}$\ldots
\note{la page s'arrête ici ; le reste de la feuille est blanc.}

\section{III. Épinglages de Legendre des courbes elliptiques}
\note{titre porté seul, au crayon, sur la page~69 de l'archive, feuillet de
couleur qui sert de page de garde à la partie ; le chiffre romain est de
l'auteur.}

\subsection{Digression. Épinglages de revêtements par rev.\ Kummériens des espaces tangents (cas relatif)}
\note{titre souligné en tête de la page~70, lu d'après le manuscrit : « Épinglages de
revêtements par \struck{\ill{} fibrés tangents} rev.\ Kummériens des
\struck{fibrés} espaces tangents (cas relatif) ».}

\page{70}
Soient
\[
  (1) \qquad f : X \longrightarrow S
\]
un morphisme lisse de dim relative $1$, $t$ une section de $X$ sur $S$,
$X^* = X \setminus t(S)$. On considère des schémas \add{$X'$,} finis sur
\struck{$X$} $X$, lisses sur $S$ de dim relative $1$, munis d'une section
$t'$ au-dessus de $t$, et « modérément ramifiés » en $t'$ sur $X$, avec un
degré de ramification $n$ --- ce qui signifie que loc.\ pour la top.\ étale
sur $X$, $X'$, $X'$ est isomorphe au voisinage de tout point $x'$ de
$t'(S)$ \struck{$t(S)$} à un $X$-schéma Kummérien de $X$ au voisinage de
$x = p(x')$, de la forme
\[
  X'_0 = \operatorname{Spec} \underline{\mathcal{O}}_X[z]/(z^n - \varphi)
\]
où $\varphi$ est une équation locale du diviseur relatif $t(S) \subset X$ au
voisinage de $x$.
\marginal{\emph{NB} $n$ est premier aux car.\ résiduelles de $S$ !!!}
\marginal{$p : X' \to X$}

Sous ces conditions, via \uncertain{l'hom.\ norme} locale, on définit un
isom.\ canonique
\[
  (2) \qquad T_{X'/S, t'}^{\otimes n} \xrightarrow[\sim]{\ \kappa_{t'}\ } T_{X/S, t}
\]
\note{dans (2), l'exposant $\otimes n$ est écrit au-dessus de la flèche, près
de $T_{X'/S,t'}$ ; dans l'indice, un $S$ corrige une première lettre.}

\page{71}
Si d'autre part $\underline{G}$ est un schéma en groupes fini étale sur $S$,
et si $X'$ est muni d'une structure de $\underline{G}_X$-torseur au-dessus
de $X^*$, alors le stabilisateur dans \struck{$\underline{G}_X$}
$\underline{G}$ de la section $t'$ est un sous-schéma en groupes
$\underline{G}_{t'}$ de $\underline{G}$, \struck{qui \ill{}} \add{qui}
\add{\uncertain{opère} \uncertain{fid.}\ sur} \struck{\ill{}}
$T_{X'/S, t'}$ et dont l'image dans
\struck{$T_{X'/S,t'}$} $\mathbb{G}_{m,S} \simeq
\underline{\mathrm{Aut}}_{\mathcal{O}_S} T_{X'/S, t'}$ est $\mu_{n,S}$,
d'où un monomorphisme canonique
\[
  (3) \qquad \mu_{n,S} \hookrightarrow \underline{G} .
\]
\marginal{\emph{NB} Cela signifie que $\underline{G}$ opère sur
$X'^* = X' | X^*$ au-dessus de $X^*$, de façon à faire de $X'^*$ un
$\underline{G}_{X^*}$-torseur sur $X^*$}
\note{la note marginale est écrite de biais dans la marge gauche, à hauteur
des premières lignes ; « au-dessus de $X^*$ » et « torseur » y sont lus d'après
le contexte. Le passage « \uncertain{opère} \uncertain{fid.}\ sur » remplace
deux mots biffés soulignés, peu lisibles.}

Par ailleurs, désignant par
\[
  f^\circ : X^* \longrightarrow S
\]
le morphisme induit par $f$, le $\underline{G}_{X^*}$-torseur $X'^*$ définit
une section
\[
  (4) \qquad \xi_{X'/S} \in H^0\bigl(S, R^1 f^\circ_* f^{\circ *}(\underline{G})\bigr) .
\]

Ceci posé, on a :

\emph{Proposition} $(X', t')$ au-dessus de $(X, t)$, avec l'opération
de $\underline{G}$ dessus, est défini à isom.\ unique près par la
connaissance de (3) (4) et de (2). \struck{Ce dernier} De façon précise,

\page{72}
\struck{pour une section} considérons (pour $S, X, t$ fixés) la catégorie des
$(X', t', \underline{G})$, \struck{où $X'$ est un $X$-schéma} \add{tels que
dessus \ill{}}, comme morphismes
\[
  (X', t', G) \longrightarrow (X'_1, t'_1, G_1)
\]
\struck{\ill{}} \add{\ill{}} des morphismes $X' \to X'_1$, et morphismes
$G \to G_1$, compatibles avec les sections marquées et les opérations.
\marginal{\emph{NB} On vérifie que pour $X', t', G$ fixés, et
$(X'_1, t'_1, G_1)$ aussi, et pour \ill{} \ill{} \ill{}, il y a au plus un
$G \to G_1$ \ill{}}
\note{la note marginale, écrite de biais dans la marge gauche, n'est lue
qu'en partie.}
Considérons d'autre part les \struck{triples} systèmes
\[
  (\underline{G}, n, i, \xi, T, k)
\]
\note{dans le sextuple, deux lettres (sans doute $\xi$ et un symbole
noirci) sont biffées et réécrites.}
où $\underline{G}$ \uncertain{un schéma en groupes fini étale sur $S$,}
$n \in \mathbb{N}^*$, $\xi \in H^0(S, R^1 f^\circ_* f^{\circ *}(\underline{G}))$
tels que les rev.\ \uncertain{correspondants} (étales sur $S$) de $X$
\uncertain{qu'elle} définit soient modérément ramifiés d'ordre $n$ en le long de
$t$, $i : \mu_{n,S} \hookrightarrow \underline{G}$ est un monomorphisme de
groupes, $T$ un Module inversible sur $S$, et $k$ un isom
\[
  T^{\otimes n} \xrightarrow{\ \sim\ } T_{X/S, t}
\]
avec la définition ad hoc des morphismes (\emph{NB} bien sûr
$\underline{G} \to \underline{G}_1$ doit induire un \struck{\ill{}}
morphisme $\mu_{n,S} \to \mu_{n_1,S}$, qui soit l'homom.\
$\lambda \mapsto \lambda^{n/n_1}$\ldots). On a un foncteur canonique de la
première catégorie dans la

\page{73}
seconde, et je dis que ce foncteur est une équivalence de catégories.

Le cas qui m'intéresse est celui où
\[
  \underline{G} = \mu_2 \underset{\text{can}}{\simeq} \{\pm 1\}_S , \quad \text{et où } n = 2
\]
--- donc où on s'intéresse aux rev.\ quadratiques $X'$ de $X/S$, étales
au-dessus de $X^*$ et « totalement ramifiés le long de la section $t$ », où
$S$ a car.\ rés.\ $\neq 2$. \struck{Ici} Pour une section
$\xi \in H^0(S, R^1 f^\circ_* (\mu_{2,X^*}))$ fixée, satisfaisant à la
condition de ramification totale, la catégorie des
\struck{revêtements} \add{\uncertain{correspondants} (\uncertain{i.e.}\ \ill{})}
\add{\ill{}} équivaut à celle des \struck{\ill{}} couples $(T, k)$ d'un
Module inv.\ $T$ sur $S$, et d'un isom
\[
  T^{\otimes 2} \simeq T_{X, t} .
\]

Le cas qui nous intéresse pour les courbes elliptiques est celui où $X$ est
de la forme $\Sigma_{\underline{J}}^*$ (fibré en droites projectives
\uncertain{marqué} \ill{} étale cubique $\underline{J}$ de $S$), associé à
un rev.\ \ill{} déduit d'une courbe elliptique $X'$ comme $X'/\{\pm 1\}$, et
où $t$ est la section de $\Sigma_{\underline{J}}^*/S$ \struck{\ill{}}
image de la section nulle de $X' = E$ \ldots
\note{les deux dernières lignes sont écrites serrées, l'une au-dessus de
l'autre, et ne sont lues qu'en partie ; « $\Sigma_{\underline{J}}^*$ » après
« de la forme » est bien marqué d'un astérisque.}

\page{74}\note{p.~1 de l'auteur, chiffrée en tête de page : une nouvelle
suite commence ici, dont la numérotation des formules repart de (1).}
Soit $E$ une courbe elliptique --- on travaille sur un corps alg.\ clos $k$
pour simplifier le propos, mais la seule hypothèse essentielle est que le
\uncertain{base} soit : car.\ résiduelles $\neq 2$. On désigne par
\struck{\ill{}}
\[
  (1) \qquad J = {}_2E(k)^*
\]
l'ens.\ des pts de $E$ d'ordre $2$, on sait que l'on a
\[
  (2) \qquad E/(\pm \mathrm{id}_E) \simeq \Sigma_J
\]
\struck{\ill{}} droite projective $J$-marquée.
\[
  (3) \qquad \Sigma_J \simeq \mathbb{P}(V_J(k)) \qquad
  0 \to V_J(k) \to k^J \xrightarrow{\ \text{somme}\ } k \to 0
\]
La donnée de $E$ est équivalente à celle de $J$ i.e.\ de $\Sigma_J$ muni du
diviseur étale $D_J$ de degré $3$, plus \struck{le} \uncertain{les} point
\[
  (4) \qquad t \in \Sigma_J^* = \Sigma_J \setminus D_J
\]
image de l'élément $0_E$, plus une droite $T$ \add{sur $k$} (savoir
$T_{E,0}$), munie d'un isomorphisme
\[
  (5) \qquad T^{\otimes 2} \xrightarrow{\ \sim\ } T_{\Sigma, t}
\]
(\struck{isomorphismes} \add{équivalence} de catégories-groupoïdes \ldots\
valable en fait sur une base $S$ quelconque ; car.\ rés.\ $\neq 2$, et
compatible aux changements de base\ldots).

Notons qu'on a un isomorphisme canonique
\[
  (6) \qquad T_\Sigma \simeq \underline{\mathcal{O}}_\Sigma(2)(\underline{\omega})
\]
où $\underline{\omega} = \underline{\omega}(J)$ est l'ens.\ des deux ordres
circulaires sur $J$, et où $L(\underline{\omega})$ (pour un objet tel que
faisceau inversible etc.)\ désigne $L \wedge_{\{\pm 1\}} \underline{\omega}$.
\note{le « 1 » en marge gauche devant « Notons » est de l'auteur ; son rôle
n'est pas clair.}

\page{75}\note{p.~2 de l'auteur.}
On a un candidat naturel (pour $\Sigma = \Sigma_J$, $D = D_J$,
$t \in \Sigma^*$ donnés) comme droite $T$ munie de (5), savoir
\[
  T_0 \simeq \mathcal{O}_t(1)
\]
à cela près qu'il y a \emph{deux} isom.\ naturels
\[
  T_0^{\otimes 2} = \mathcal{O}_t(2)
  \xrightarrow{\ \sim\ } T_{\Sigma, t} \simeq \mathcal{O}_t(2)(\underline{\omega}) ,
\]
correspondant aux deux éléments $\omega$ de $\underline{\omega}$, savoir
\[
  (7) \qquad \xi \longmapsto \xi \wedge \omega \qquad
  T_0^{\otimes 2} \xrightarrow{\ \sim\ } \mathcal{O}_t(2) \simeq T_{\Sigma, t}
  = \underline{\mathcal{O}}_t(2)(\underline{\omega})
\]
\note{les numéros (7) et (8) des deux premières formules, \struck{(7)} et \struck{(8)}, sont biffés ; la
troisième reçoit le numéro (7). Devant $\xi \mapsto \xi \wedge \omega$, un
signe biffé. En tête de page, $\Sigma_{\underline{J}}$, $D_{\underline{J}}$
sont d'abord écrits puis surchargés, les égalités $= \Sigma_J$, $D = D_J$ étant
ajoutées au-dessus.}

On va lever cette ambiguïté, en remplaçant $T_0 = \mathcal{O}_t(1)$ par une
variante tordue. Le tordre par un torseur sous $\mu_2$ \add{$= \mu_2(k)$}
\ill{} ne \ill{} pas, car $T_0^{\otimes 2}$ n'\ill{} \ill{} pas changé
(à isom.\ can.\ près) --- on le tordra par un torseur (qui ne précise)
$Q_0$ sous $\mu_4 = \mu_4(k)$ (qui peut être \uncertain{interprété} comme
l'ens.\ des 4 sommets d'un \emph{carré}, dont $\mu_4$ est le groupe des
rotations, \ill{} $\{i, -i\} = \mu_4^*$ \ill{} \ill{} \struck{\ill{}}
\add{\ill{}} n'y \uncertain{jouent} \uncertain{pas} de \uncertain{rôle}
privilégié\ldots). On pose donc
\[
  (8) \qquad T_0 = \underline{\mathcal{O}}_t(1) \wedge_{\mu_4} Q_0
\]
d'où
\[
  (9) \qquad T_0^{\otimes 2} \simeq \underline{\mathcal{O}}_t(2) \wedge_{\mu_2} D_{Q_0}
\]
où
\[
  (10) \qquad D_{Q_0} = Q_0/\{\pm 1\}
\]
est l'un des deux diagonales du carré $Q_0$. Supposons donc qu'on se donne
un carré \add{combinatoire (« 4 sommets »),} $Q_0$, avec plus
\uncertain{isomorphisme}. \struck{\ill{} $Q_0$ \ill{} \ill{} le groupe des}
\[
  (11) \quad
  \begin{cases}
    \underline{\omega}_{Q_0} \ (\text{ens.\ des orientations de } Q_0)
      \simeq \mu_4^* = \{i, -i\} \\
    D_{Q_0} \ (\text{ens.\ des diagonales de } Q_0) \simeq \underline{\omega}
  \end{cases}
\]
\note{dans (10), « l'un des deux diagonales » est la lecture de la page ;
le sens attendu est « l'ens.\ des deux diagonales ».}

\page{76}
Alors, définissant $T_0$ par (8) (pour $t \in \Sigma_J^*$ donné) on a un
isom.\ canonique
\[
  (12) \qquad T_0^{\otimes 2} \xrightarrow{\ \sim\ } T_{\Sigma, t}
\]
Ceci posé, on appellera \struck{\ill{}} \emph{épinglage de Legendre}
de $E$, subordonné au carré $Q_0$ (muni des données (11)) \struck{\ill{}}
\add{un} isomorphisme
\[
  (13) \qquad T = T_{E, 0} \xleftarrow[\ \sim\ ]{\ \kappa\ } T_0
\]
compatible avec les \uncertain{isom.\ normes}, i.e.\ rendant commutatif le
diagramme
\note{avant « de $E$ », deux lignes biffées, illisibles ; « épinglage » y
remplace un mot biffé.}


\begin{tikzcd}
  T_0^{\otimes 2} \arrow[rr, "\kappa^{\otimes 2}", "\sim"'] \arrow[dr, "(12)"'] & & T^{\otimes 2} \arrow[dl, "(5)"] \\
  & T_{\Sigma, t} &
\end{tikzcd}

(14).

Il y en a exactement deux, différant par le signe, i.e.\ formant un torseur
sous $\mu_2$.

Bien entendu, on veut pouvoir définir la notion d'épinglage de $E$
(\uncertain{en termes} d'un $Q_0$ donné \uncertain{une fois pour toutes}) pour
$E$ variable --- mais dans ce cas $\underline{\omega} = \underline{\omega}_J$
n'est pas encore déterminé (puisque $J$ ne l'est pas). Donc la donnée d'un
épinglage $(\alpha, \kappa)$ de Legendre sera définie comme le couple formé
d'un isom.
\[
  (15) \qquad \alpha : \underline{\omega}_J \xrightarrow{\ \sim\ } D_{Q_0} = Q_0/\{\pm 1\}
\]
\emph{et} d'un isom.\ $\kappa$ (13) satisfaisant la compatibilité (14),
\note{dans (14), un trait vertical est esquissé sous $T_0^{\otimes 2}$ à
côté de la flèche marquée (12) ; il n'est pas repris.}

\page{77}\note{p.~3 de l'auteur.}
où $Q_0$ est un carré donné une fois pour toutes (la chose fixée, disons,
quand même), avec comme seule donnée un iso
\[
  (16) \qquad \underline{\omega}_{Q_0} \simeq \mu_4^*
  \qquad \text{i.e.\ d'un élément de } \mu_4^* \wedge \underline{\omega}_{Q_0}
\]
i.e.\ un isom.
\[
  (17) \qquad \mu_4^* \simeq \operatorname{Aut}^+(Q_0) .
\]
\note{dans (16), une lettre noircie en indice de $\underline{\omega}$ est
réécrite $Q_0$ au-dessous ; le numéro (17) est surchargé.}

Pour des raisons qui apparaîtront par la suite, le choix le plus
\uncertain{commode} est de prendre pour $Q_0$ le carré \struck{\ill{}
satisfaisant} \add{\uncertain{i.e.\ donné}} (17), avec un côté fixé, dont
\struck{les \ill{}} \add{l'un des \uncertain{sommets} \uncertain{incidents}}
$s, s'$ \struck{\ill{}} dans un \uncertain{sens} \uncertain{convenable}
\ill{} avec l'un $\mu_4^* = \{i, i'\}$ des racines carrées de $-1$, par la
condition que
\[
  (18) \qquad s_{i'} = i\, s_i \qquad \text{d'où} \qquad s_i = i' s_{i'}
\]
puisque $i' = i^{-1} = -i$. Ainsi, $T_0$ peut s'interpréter comme munie de
la structure \struck{\ill{}} \add{suivante, qui la caractérise :}
$T_0 \otimes \underline{\mathcal{O}}_t(-1)$ est muni d'une paire de base,
\add{en corr.\ 1-1} \struck{\uncertain{indexée} \ill{}} avec $\mu_4^*$,
satisfaisant les deux relations (18),\ i.e.\ maintenant le produit $i s_i$
est \ill{} dans la structure vectorielle de
$T_0 \otimes \underline{\mathcal{O}}_t(-1)$. \struck{\ill{}}

Ici on aura un iso.\ canonique
\[
  (19) \qquad D_{Q_0} \simeq \mu_4^*
\]
et par suite la donnée (16) équivaut à la donnée d'un iso
\note{en marge gauche, à hauteur de (18), un petit carré dont le sommet
supérieur droit est marqué $s' = s_{i'}$ et le sommet inférieur droit
$s = s_i$ ; le côté droit est le côté fixé. Dans la marge de (18), devant
$s_{i'}$, une égalité biffée noircie ; une autre esquisse biffée
($s' = is$ ?) se trouve à la ligne au-dessus, et en tête de (18) une
égalité biffée \struck{$s' = is$}.}

\page{78}
\[
  (20) \qquad \alpha : \underline{\omega}_J \simeq \mu_4^*
  \qquad \text{i.e.\ d'un élément de } \underline{\omega}_J \wedge \mu_4^* .
\]
Ces définitions posées, on voit donc que la catégorie\add{-groupoïde} des
courbes elliptiques \add{$/k$} épinglées de Legendre \struck{\ill{}}
\add{est} équivalente à la catégorie-groupoïde des \struck{triplets
couples} \add{triplets} $(J, t, \alpha)$ \struck{de} formés
\[
  (21) \quad
  \begin{cases}
    \text{a) d'un $J$ à trois éléments (d'où $\Sigma_J$, $D_J$, $\Sigma_J^*$),} \\
    \text{b) d'un $t \in \Sigma_J^*(k)$} \\
    \text{c) d'un isom $\alpha : \underline{\omega}_J \simeq \mu_4^*$ \ \ $(= \mu_4^*(k))$}
  \end{cases}
\]

Un épinglage de Legendre--Jacobi d'échelon 2 \add{d'une courbe elliptique
$E$} sera la donnée d'un couple d'épinglages, l'un de Legendre, l'autre de
Jacobi d'échelon 2 i.e.\ d'une bijection $J \simeq \{0, 1, \infty\}$,
permettant d'identifier $\Sigma_J$ à la droite projective type
$\mathbb{P}^1_k$. On a alors dans
$\underline{\omega}_J \simeq \underline{\omega}_{\{0,1,\infty\}}$ un élément
\struck{\ill{}} privilégié $\omega^+$, et la donnée c) de l'épinglage de
Legendre équivaut \add{donc} à celle d'une racine carrée
\struck{\ill{}} de $-1$, soit $i$ --- c'est la donnée « arithmétique ».
\struck{Par ailleurs, la catég.}\ La catégorie-groupoïde des courbes
elliptiques sur $k$, munies d'un épinglage de Legendre--Jacobi,
\add{est rigide (\uncertain{pas} d'auto.\ \struck{\ill{}})}
\struck{équivalente à} \struck{celle des \ill{} à isom.\ près ce} et elle
\note{« Legendre » dans « d'un épinglage de Legendre » est d'abord écrit
avec un « à » d'une autre forme ; dans c), un symbole biffé suit $\mu_4^*$.
La dernière ligne, biffée, et l'ajout au-dessus (« est rigide \ldots »),
ne sont lus qu'en partie. Dans a), avant $J$, le mot \struck{ens.}\ est biffé.}

\page{79}\note{p.~4 de l'auteur.}
est équivalente à celle des couples $\{t, i\}$, avec
\[
  t \in U_{0,3}(k) = \mathbb{P}^1_k(k) \setminus \{0, 1, \infty\}, \quad \text{et } i \in \mu_4^* .
\]

\struck{On peut définir sur $\Sigma_J$ le fibré
$\mathcal{O}(1)^{\natural}_{\Sigma_J} = \underline{\mathcal{O}}_{\Sigma_J}(1) \wedge_{\mu_4}$}
\note{la phrase et la formule sont biffées d'un trait ; un trait ondulé
relie la formule biffée à la ligne suivante. L'exposant écrit ici $\natural$
est un signe en zigzag, proche de « ϟ », propre à cette page et à la
suivante ; on le rend partout par $\natural$. Il est distinct du $\flat$ de
(23).}

J'opère : sur une base quelconque, mais \uncertain{éventuellement} en car.\
rés.\ $\neq 2$ i.e.\ sur la base absolue \add{$S_0 =$}
$\operatorname{Spec}(\mathbb{Z}[1/2])$. Sur \struck{$\Sigma_S$} $S_0$ je
définis, l'opération de \uncertain{torsion} $\natural$, sur les
\struck{modules} faisceaux en Modules, \struck{par} sur un schéma de base
$S$
\[
  (22) \qquad M^{\natural} = M \wedge_{\mu_{4S}} Q_{0S}
  \qquad \underline{\text{NB}}\ \ M^{\natural} \simeq M \otimes_{\mathcal{O}_S} \underline{\mathcal{O}}_S^{\natural}
\]
où ici on considère $Q_0$ comme un $\mu_4$-torseur sur la base absolue
$S_0$. On aura donc
\[
  (23) \qquad M^{\natural} \otimes M^{\natural} \simeq M^{\flat}
  \overset{\text{déf}}{=} M \wedge_{\mu_{2S} = \{\pm 1\}_S} \mu_{4S}^*
\]
où \struck{$\mu_4^* \simeq \ldots$}
\[
  (24) \qquad \mu_4^* = \mu_4 \setminus \mu_2
\]
est considéré comme torseur sous $\mu_2$, i.e.\ rev.\ étale quadratique
de la base absolue $S_0$.

On s'intéressera donc, \struck{\ill{}} pour un \uncertain{revêtement}
\note{sous $Q_{0S}$ dans (22), un petit mot biffé ; la dernière ligne se
poursuit page~80.}

\page{80}
étale $\underline{J}$ de degré $3$ \struck{de} de $S$, permettant de
définir un $\Sigma_{\underline{J}}$
\[
  (25) \qquad \Sigma_{\underline{J}} = \mathbb{P}(V_{\underline{J}}(\underline{\mathcal{O}}_S))
  \qquad 0 \to V_{\underline{J}}(\mathcal{O}_S) \to
  \underset{\substack{\parallel \\ p_*(\mathcal{O}_{\underline{J}}) \\ p : \underline{J} \to S}}{\underline{\mathcal{O}}_S^{\underline{J}}}
  \xrightarrow{\ \mathrm{tr}\ } \underline{\mathcal{O}}_S
\]
et au module inversible $\underline{\mathcal{O}}_{\Sigma_{\underline{J}}}(1)$,
au module
\[
  (26) \qquad T_{0, \Sigma_{\underline{J}}} = \underline{\mathcal{O}}_{\Sigma_{\underline{J}}}(1)^{\natural}
\]
de sorte qu'on a un iso
\[
  (27) \qquad T_{0, \Sigma}^{\otimes 2} \simeq \underline{\mathcal{O}}_{\Sigma_J}(2) \wedge_{\{\pm 1\}} \mu_4^*
  \simeq T_{\Sigma_J} \wedge_{\{\pm 1\}} (\mu_4^* \wedge \underline{\omega}_{\underline{J}})
\]
où $\underline{\omega}_J$ est le revêtement \add{étale} quadratique associé au
revêtement \add{étale} cubique $\underline{J}$.

Si maintenant $E$ est une courbe elliptique sur la base $S$, on lui
associe
\[
  (28) \qquad \underline{J} = ({}_2E)^* \qquad \text{rev.\ étale cubique de } S
\]
et on \struck{\ill{}} a un iso
\[
  (29) \qquad E/(\pm \mathrm{id}_E) \xrightarrow{\ \sim\ } \Sigma_{\underline{J}}
\]
et un iso
\note{dans (25), la trace $\mathrm{tr}$ est écrite au-dessus de la flèche,
d'une lecture incertaine ; l'identification
$\underline{\mathcal{O}}_S^{\underline{J}} = p_*(\mathcal{O}_{\underline{J}})$,
$p : \underline{J} \to S$, est portée au-dessous. Le « au module inversible
\ldots, au module » est lu d'après le contexte. La page s'arrête sur « et un
iso », la suite dans le lot suivant.}

\end{document}
